% =================================================
% Ученический рабочий лист
% =================================================

\worksheettitle
  {Длина отрезка}
  {Измерение, единицы длины и построение}

\studentnameblock

\begin{checkpointbox}[Входная разминка]
  \begin{enumerate}[label=\textbf{\arabic*.},itemsep=0.45em]
    \item Какой геометрический объект можно измерить: прямую, луч или
      отрезок? \answerblank[38mm]
    \item Каким инструментом измеряют и строят отрезки?
      \answerblank[42mm]
  \end{enumerate}
\end{checkpointbox}

\section{Отрезок и его длина}

\segmentlengthfigure

\begin{fillbox}[Дополните]
  Отрезок — это геометрическая \answerblank[30mm], а его длина —
  \answerblank[28mm] с единицей измерения.

  Запись \(AB=6\text{ см }4\text{ мм}\) сообщает
  \answerblank[50mm] отрезка \(AB\).
\end{fillbox}

\begin{practicebox}[Задание 1. Верно или неверно?]
  Поставьте знак «\(+\)» или «\(-\)».
  \begin{enumerate}[label=\textbf{\arabic*.},itemsep=0.38em]
    \item \choicebox\ Отрезок и его длина — одно и то же.
    \item \choicebox\ Длина отрезка записывается вместе с единицей измерения.
    \item \choicebox\ Длину всей прямой можно измерить линейкой.
    \item \choicebox\ Равные отрезки имеют одинаковую длину.
  \end{enumerate}
\end{practicebox}

\section{Измеряем отрезок}

\correctmeasurementfigure

\begin{fillbox}[Алгоритм]
  Один конец отрезка совмещают с делением \answerblank[14mm] линейки.
  Затем определяют, с каким делением совпадает \answerblank[35mm] конец.
\end{fillbox}

\readrulerpracticefigure

\begin{practicebox}[Задание 2. Прочитайте показания линейки]
  \begin{enumerate}[label=\alph*),itemsep=0.45em]
    \item \(AB=\answerblank[40mm]\);
    \item \(CD=\answerblank[40mm]\).
  \end{enumerate}
\end{practicebox}

\begin{thinkbox}[Задание 3. Отрезок начинается не у нуля]
  На нижнем рисунке точка \(C\) находится у деления \(1\text{ см }3\text{ мм}\),
  а точка \(D\) — у деления \(7\text{ см }1\text{ мм}\).
  Запишите вычисление длины \(CD\):
  \[
    CD=\answerblank[28mm]-\answerblank[28mm]=\answerblank[34mm].
  \]
\end{thinkbox}

\section{Единицы длины}

\unitsladderfigure

\begin{fillbox}[Заполните равенства]
  \[
    1\text{ см}=\answerblank[16mm]\text{ мм},\qquad
    1\text{ дм}=\answerblank[16mm]\text{ см},
  \]
  \[
    1\text{ м}=\answerblank[16mm]\text{ см},\qquad
    1\text{ км}=\answerblank[20mm]\text{ м}.
  \]
\end{fillbox}

\begin{practicebox}[Задание 4. Выразите в указанных единицах]
  \begin{enumerate}[label=\alph*),itemsep=0.5em]
    \item \(7\text{ см }9\text{ мм}=\answerblank[25mm]\text{ мм}\);
    \item \(5\text{ м }8\text{ см}=\answerblank[25mm]\text{ см}\);
    \item \(905\text{ см}=\answerblank[22mm]\text{ м }\answerblank[18mm]\text{ см}\);
    \item \(6\text{ км }45\text{ м}=\answerblank[28mm]\text{ м}\).
  \end{enumerate}
\end{practicebox}

\begin{warningbox}[Проверьте себя]
  Ученик написал:
  \[
    6\text{ км }45\text{ м}=645\text{ м}.
  \]
  Объясните ошибку и запишите верный ответ.
  \writinglines{1}
\end{warningbox}

\section{Строим отрезки}

\begin{fillbox}[Расставьте шаги по порядку]
  Запишите номера от \(1\) до \(4\).
  \begin{enumerate}[label=\choicebox,itemsep=0.35em]
    \item отметить второй конец у нужного деления;
    \item соединить отмеченные точки;
    \item отметить первый конец;
    \item совместить с ним деление \(0\) линейки.
  \end{enumerate}
\end{fillbox}

\begin{practicebox}[Задание 5. Выполните построения линейкой]
  \studentconstructionfigure
\end{practicebox}

\section{Равные отрезки}

\equalcontrastingsegmentsfigure

\begin{practicebox}[Задание 6]
  Какие отрезки равны? \answerblank[48mm]

  Объяснение: \answerblank[95mm]
\end{practicebox}

\section{Точка внутри отрезка}

\partssegmentfigure

\begin{fillbox}[Главное равенство]
  Если точка \(C\) лежит между \(A\) и \(B\), то
  \[
    AB=\answerblank[22mm]+\answerblank[22mm].
  \]
\end{fillbox}

\partsproblemfigure

\begin{practicebox}[Задание 7]
  Найдите длину \(CB\).
  \[
    CB=\answerblank[30mm]-\answerblank[30mm]=\answerblank[35mm].
  \]
\end{practicebox}

\begin{checkpointbox}[Задание 8. Самостоятельно]
  Точка \(K\) лежит между точками \(M\) и \(N\).
  \(MN=12\text{ см}\), \(MK=7\text{ см }6\text{ мм}\).
  Найдите \(KN\).
  \writinglines{2}
\end{checkpointbox}

\begin{reflectionbox}[Итог]
  Объясните, почему по одному только рисунку нельзя надёжно определить
  длину отрезка, если не дан масштаб и не выполнено измерение.
  \writinglines{1}
\end{reflectionbox}
